\subsubsection{LLVM XRay}\label{llvm-xray}

\href{https://llvm.org/docs/XRay.html}{LLVM XRay} je instrumentacija koju Clang može da ubaci u funkcije programa. Za razliku od \texttt{gprof}-a, XRay je projektovan za selektivno uključivanje i detaljnije tragove ulaska i izlaska iz funkcija. Koristan je kada želimo da vidimo vremenski trag poziva, latencije funkcija ili kada želimo da izdvojimo skupe putanje kroz program.

Primer ima tri tipa posla: CPU petlju, sortiranje i povremeno čekanje:

\codefile{xray_workload.cpp}{cpp}

Prevodimo ga Clang-om sa XRay instrumentacijom:

\begin{minted}{bash}
$ make
\end{minted}

XRay se uključuje pri pokretanju kroz promenljivu okruženja:

\begin{minted}{bash}
$ XRAY_OPTIONS="patch_premain=true xray_mode=xray-basic verbosity=1" ./xray_workload
\end{minted}

Pokretanje pravi datoteku oblika \texttt{xray-log.xray\_workload.\textless{}pid\textgreater{}}. Tekstualni sažetak se može dobiti alatom \texttt{llvm-xray}:

\begin{minted}{bash}
$ llvm-xray account xray-log.xray_workload.* > xray-account.txt
$ less xray-account.txt
\end{minted}

U sažetku tražimo funkcije sa najvećim ukupnim ili prosečnim vremenom. U ovom primeru \texttt{wait\_work} ima veliko zidno vreme po pozivu, dok \texttt{busy\_work} i \texttt{sort\_work} pokazuju aktivan rad programa.

Za trag poziva možemo koristiti \texttt{stack} analizu:

\begin{minted}{bash}
$ llvm-xray stack xray-log.xray_workload.* > xray-stacks.txt
\end{minted}

\paragraph{Kada koristiti XRay}\label{xray-kada}

XRay je dobar kada je pitanje vezano za funkcijske pozive i latenciju putanja kroz program: koji handler je spor, koliko traje određena faza, koliko često se ulazi u skupu funkciju. Pošto je instrumentacija deo binarne verzije, treba praviti posebnu konfiguraciju za merenje i jasno je odvojiti od produkcione verzije.

\paragraph{Ograničenja}\label{xray-ogranicenja}

\begin{itemize}
\tightlist
\item
  Potrebni su Clang, XRay runtime i \texttt{llvm-xray}; dostupnost zavisi od distribucije i arhitekture.
\item
  Instrumentacija menja binarnu datoteku i uvodi dodatni trošak.
\item
  Veoma česte male funkcije mogu napraviti veliki log. U praksi treba ograničiti šta se instrumentira ili koristiti pragove.
\item
  XRay meri funkcijske događaje; za hardverske događaje, keš i mikroarhitekturu treba koristiti alate kao što su \texttt{perf} ili VTune.
\end{itemize}
